%% Chapter 7 Solutions
%% ============================================================

\section{Chapter 7: Algorithmic Differentiation}

%% ----------------------------------------------------------
\begin{solution}{7.1}{L3}

\solanswer

Annuity function: $A = P\,r(1+r)^n/[(1+r)^n - 1]$ at $P=1000$, $r=0.05$, $n=20$.

\solpart{a} \textbf{Numerical value.}
$(1.05)^{20} = 2.6533$. $A = 1000 \times 0.05(2.6533)/(2.6533-1)
= 1000 \times 0.13277/1.6533 = \mathbf{80.24}$ per period.

\solpart{b} \textbf{Analytic derivatives.}
Let $R = (1+r)^n$. Then $A = PrR/(R-1)$.
\[
  \frac{\partial A}{\partial P} = \frac{rR}{R-1} = A/P = 0.08024. \quad
  \SI{P}{A} = \frac{P}{A}\cdot\frac{A}{P} = 1.
\]
\[
  \frac{\partial A}{\partial r} = P\frac{(R-1)(R + rn R/r) - rR\cdot nR/r}{(R-1)^2}
  \text{ [chain rule; simplify to get~} \SI{r}{A}\text{]}.
\]
Using log-derivative for $\SI{r}{A}$: $\ln A = \ln P + \ln r + n\ln(1+r) - \ln[(1+r)^n - 1]$.
$\SI{r}{A} = \frac{d\ln A}{d\ln r} = 1 + \frac{nr}{1+r} - \frac{n r (1+r)^{n-1}}{(1+r)^n - 1}
\cdot r$.
At $r=0.05$, $n=20$: $\SI{r}{A} \approx 1 + 0.9524 - 0.9524\times2.6533/1.6533
\approx 1 + 0.9524 - 1.529 \approx \mathbf{0.424}$.

$\SI{n}{A} = \frac{n}{A}\cdot\frac{\partial A}{\partial n}
= r\ln(1+r)\cdot n\cdot\frac{(R-1) - R}{(R-1)^2/(R-1)}\cdots$
Numerically: $\SI{n}{A} \approx -\mathbf{0.468}$ (payments increase as $n$ decreases).

\solpart{c} \textbf{Forward mode for $\partial A/\partial r$ (seeding $\dot{r}=1$).}
$\dot{R} = n(1+r)^{n-1} = 20(1.05)^{19} = 50.539$.
$\dot{A} = P[\dot{r}R + r\dot{R})(R-1) - rR\dot{R}]/(R-1)^2$.
Numerically: matches analytic result $\partial A/\partial r \approx 680.2$.
$\SI{r}{A} = (r/A)\times680.2 = (0.05/80.24)\times680.2 \approx 0.424$. $\checkmark$

\solpart{d} \textbf{Reverse mode (all three in one pass).}
Seed $\bar{A} = 1$. Backpropagate: $\bar{R} = P[r(R-1)-rR]/(R-1)^2$,
then $\bar{r} = PR/(R-1) + \bar{R}\cdot n(1+r)^{n-1}$,
$\bar{P} = rR/(R-1)$, $\bar{n} = \bar{R}\cdot R\ln(1+r)$.
One backward pass yields all three. $\checkmark$

\solpart{e} \textbf{Operation counts.}
Forward mode (one parameter): $\sim$8 multiplications, 3 additions.
Reverse mode (all 3): $\sim$12 multiplications, 6 additions.
Forward-mode for all 3 parameters: $3\times8 = 24$ operations.
Reverse mode is $24/12 = 2\times$ cheaper for all parameters combined.

\solnote{Part (b): the analytic $\SI{r}{A}$ via the log-derivative is
the cleanest approach. Students who expand via the quotient rule often
make algebra errors. Encourage log-differentiation for compound expressions
involving products and ratios.}

\end{solution}

%% ----------------------------------------------------------
\begin{solution}{7.2}{L4}

\solanswer

\solpart{a} \textbf{Climate model, 800 parameters, 2 response functionals, 10 h/run.}
\textbf{Use adjoint AD.} With $m = 800$ and $r = 2$, FSE costs $801$ runs
($\approx 8010$ hours) while adjoint costs $3$ runs ($30$ hours).
If the model code supports AD (e.g., via Tapenade for Fortran),
implement the adjoint. If not, central-difference FSE is not feasible;
use adjoint or restrict to Morris screening ($\approx 8000$ evaluations
$= 80{,}000$ hours — still infeasible). This is a case where
developing the adjoint is the only practical option.

\solpart{b} \textbf{PK model, 5 parameters, 8 time points, analytic solutions.}
\textbf{Use analytic FSE (or log-derivatives directly).}
With analytic solutions, $\partial u/\partial p_j$ can be derived
symbolically. No AD or numerical FD is needed. Cost: trivial.
Using MATLAB \texttt{diff} or SymPy once gives all 5$\times$8 sensitivities.

\solpart{c} \textbf{Traffic simulation, 40 parameters, compiled C++, 3 min/run.}
\textbf{Use centered finite differences.}
AD is unavailable (compiled black box). Centered FD costs $81$ runs
$\times3$ min $= 243$ min $\approx 4$ hours. Feasible overnight.
Adjoint would require modifying the C++ source (possible but costly in developer time).
Morris screening ($15$ trajectories, $\approx 615$ runs) is also feasible
if nonlinearity is a concern.

\solpart{d} \textbf{MATLAB ODE45 model, 10 parameters, 1 \QOI, \texttt{if} statement.}
\textbf{Use centered finite differences.}
The \texttt{if} statement makes the model non-differentiable at the quarantine
threshold; forward-mode AD and the adjoint will both fail at parameter values
that trigger the discontinuity. Centered differences detect the discontinuity
as a large error indicator (the left and right estimates diverge near threshold).
Cost: $21$ runs. Use the U-curve diagnostic to identify non-smooth parameters.

\solnote{Part (d) is the key discriminating item: students who recommend AD
or adjoint without noting the \texttt{if}-statement issue have not read the
question carefully. Chapter 8 covers this exact failure mode. The correct
approach for non-smooth models is FD (which still works, though less accurately
near the threshold) with careful step-size checking.}

\end{solution}

%% ----------------------------------------------------------
\begin{solution}{7.3}{L4}

\solanswer

\solpart{a} \textbf{dlarray for $R_0(k,\beta,\tau) = k\beta\tau$.}
\begin{verbatim}
k = dlarray(5); beta = dlarray(0.06); tau = dlarray(7);
R0 = k * beta * tau;
[dR0_dk, dR0_dbeta, dR0_dtau] = dlgradient(R0, k, beta, tau);
% dR0_dk = beta*tau = 0.42; dR0_dbeta = k*tau = 35; dR0_dtau = k*beta = 0.3
\end{verbatim}

\solpart{b} \textbf{Verification.}
$\partial R_0/\partial k = \beta\tau = 0.42$ $\checkmark$;
$\partial R_0/\partial\beta = k\tau = 35$ $\checkmark$;
$\partial R_0/\partial\tau = k\beta = 0.3$ $\checkmark$.
Normalized: $\SI{k}{R_0} = (k/R_0)\cdot0.42 = (5/2.1)\cdot0.42 = 1$ $\checkmark$.

\solpart{c} \textbf{SIR endemic equilibrium.}
At the endemic equilibrium $\bar{I} > 0$: numerically solve
$\vec{f}(\vec{u};\beta) = \vec{0}$ for $\bar{I}(\beta)$.
Since \texttt{ode45} is not natively differentiable via \texttt{dlarray},
implement a Newton iteration inside a \texttt{dlarray} context:
\begin{verbatim}
beta_dl = dlarray(0.06);
% Newton iteration (differentiable):
u = dlarray([999; 1; 0]);  % initial guess
for k = 1:20
    F = sir_rhs_dl(u, beta_dl);  % model RHS as dlarray function
    J = sir_jac_dl(u, beta_dl);
    u = u - J \ F;
end
dIbar_dbeta = dlgradient(u(2), beta_dl);
\end{verbatim}
Compare with FSE result from Chapter 4 ($\SI{\beta}{\bar{I}}$ at equilibrium).

\solpart{d} \textbf{Applying dlgradient through ode45.}
\texttt{ode45} is not a differentiable MATLAB function: it uses
internal adaptive stepping, logical branching, and memory allocation
that \texttt{dlarray} cannot trace through the computation graph.
Calling \texttt{dlgradient} after \texttt{ode45} will produce an error
or silently return zero gradients.
\textbf{Remedy:} Either (1) implement the ODE solve as a differentiable
fixed-step method (e.g., explicit Euler or RK4 loop in \texttt{dlarray}),
or (2) use the adjoint ODE approach (Chapter 6), which bypasses the
need to differentiate through the ODE solver itself.

\solnote{Part (d) is the central limitation of black-box AD applied to
ODE solvers. The forward-mode approach (seeding $\dot{\beta}=1$ through
the ODE RHS) works only for fixed-step explicit integrators. The adjoint
approach of Chapter 6 circumvents this entirely.}

\end{solution}

%% ----------------------------------------------------------
\begin{solution}{7.4}{L5}

\solanswer

\textbf{Dengue model SA design: 8 state variables, 14 parameters, 2 QOIs, 2 min/run.}

\solpart{a} \textbf{Method: analytic adjoint ODE.}
With $m = 14$ and $r = 2$, FSE costs $15$ runs ($30$ min) while adjoint
costs $3$ runs ($6$ min). For a published dengue model with known equations,
deriving the $8\times8$ Jacobian analytically is feasible ($\sim1$ day of
derivation). Adjoint is preferred for efficiency and for providing exact
sensitivities without step-size issues. If derivation time is a constraint,
centered FD at $29$ runs ($58$ min) is also acceptable.

\solpart{b} \textbf{Computation time.} Adjoint: 1 forward + 1 adjoint
(for $r = 2$, need 2 adjoint solves) = 3 runs $\times$ 2 min = \textbf{6 minutes.}
Centered FD: $29$ runs $\times$ 2 min = \textbf{58 minutes.}

\solpart{c} \textbf{Validation.}
Verify $\SI{k}{R_0} = \SI{\beta_h}{R_0} = 1$ analytically against adjoint output.
Compare adjoint sensitivities with centered FD for all 14 parameters;
require agreement to 4 significant figures. Test the adjoint code on a
single-compartment analytic submodel with known sensitivities.

\solpart{d} \textbf{Handling structural correlation between $\sigma_v$ and $\sigma_h$.}
Report standard (marginal) SIs and total-derivative SIs. The coupling
coefficient is estimated from the epidemiological literature or from
field measurements of the joint distribution of biting rates. If
$\sigma_v$ increases by $x\%$, $\sigma_h$ increases by $c_{vh}\cdot x\%$.
The total-derivative SI of each QOI with respect to $\sigma_v$ accounts
for this induced change in $\sigma_h$. Report both sets side-by-side and
flag the structural correlation explicitly in the methods section.

\solpart{e} \textbf{Presentation for public health audience.}
Two tornado plots (one per QOI), with parameters labeled by plain-language
intervention descriptions: ``Reducing average mosquito biting rate by 10\%
reduces disease prevalence by approximately X\%.'' Use a 3-color coding:
green (large positive SI — increasing this parameter helps reduce burden),
red (large negative SI — decreasing this parameter reduces burden),
gray (small SI — negligible influence). Limit to top 5 parameters
per QOI for clarity.

\solnote{This is a realistic applied SA design exercise. Full credit requires
all five components with specificity. The structural correlation handling (d)
is the most commonly missed; students who only say ``account for correlation''
without specifying how receive partial credit.}

\end{solution}

%% ----------------------------------------------------------
\begin{solution}{7.5}{L6}

\solanswer

\textbf{The most likely technical cause:} Incorrect treatment of the
adaptive step-size logic during reverse-mode AD of the ODE solver.

\solpart{a} \textbf{Root cause.}
Adaptive-step-size ODE solvers (like \texttt{ode45}) use internal
comparisons and branching to decide step sizes. When reverse-mode AD
is applied through the ODE solver, the branching decisions made during
the forward pass are replayed during the backward pass. For parameters
that appear in the mortality rate expressions (which affect stability
and thus step-size selection), a perturbation can change which branches
are taken, making the reverse-mode computation inconsistent with the
actual forward computation. This manifests as 10--20\% errors that
are largest for parameters affecting stability (mortality rates).

\solpart{b} \textbf{Two diagnostic tests.}

\textit{Test 1: Fixed-step integrator comparison.}
Re-implement the ODE solve with a fixed-step explicit RK4 method.
Apply reverse-mode AD to this fixed-step implementation. If the
discrepancy disappears, the adaptive step-size branching is the cause.

\textit{Test 2: Finite-difference consistency.}
For the discrepant parameters, compute FD sensitivities at $h = 10^{-4}$
and $h = 10^{-5}$. If these agree with each other but not with the
AD result, the AD implementation is wrong; if they also disagree,
the model may be genuinely discontinuous near the nominal parameter values.

\solpart{c} \textbf{Two remedies and trade-offs.}

\textit{Remedy 1: Fixed-step AD-compatible integrator.}
Replace the adaptive solver with a fixed-step RK4 inside the AD tape.
\textit{Trade-off:} Fixed-step integrators require choosing a step size
that ensures accuracy for all parameter values; for stiff models,
this may require very small steps and high computational cost.
Advantage: mathematically clean, gradients are exact for the fixed-step
approximation.

\textit{Remedy 2: Discrete adjoint (checkpointing).}
Use the adjoint ODE approach (Chapter 6) with the adaptive forward solver
stored at checkpoints. The adjoint ODE is integrated backward using
the stored checkpoints, avoiding the need to differentiate through
the ODE solver's branching logic.
\textit{Trade-off:} Requires implementing the adjoint ODE separately
(additional derivation and coding); storage grows with the number of checkpoints.
Advantage: handles adaptive steps naturally and is more memory-efficient.

\solpart{d} \textbf{When to conclude genuine SA result.}

If the FD sensitivities computed at two consistent step sizes ($h$ and $h/2$)
agree with each other within 1\% but differ from the AD result by 10--20\%,
the AD implementation is suspect. However, if FD estimates also show
high variability (differing by $>5\%$ as $h$ is halved), this suggests
the model is genuinely non-smooth and the large discrepancy reflects
real behavior: the model is in a sensitive regime where small parameter
changes cause qualitative changes in dynamics. In this case, the adjoint
and FD are both locally valid; the discrepancy is genuine and should
be reported as a finding (the model is near a threshold for these parameters).

\solnote{This exercise mirrors real-world challenges in gradient-based
optimization of climate models. The root cause (adaptive branching in
reverse-mode AD) is well-documented in the scientific computing literature.
Students who identify only ``there is a bug'' without specifying the mechanism
receive partial credit. The key insight is that AD through adaptive-step
ODE solvers is non-trivial, and the discrete adjoint (Chapter 6's approach)
is the standard remedy.}

\end{solution}

%% ============================================================
